This section reviews the three areas that intersect in the paper's contribution: cooperative cable-suspended aerial transport (Section~\ref{sec:rw_transport}), fault-tolerant control for UAVs (Section~\ref{sec:rw_ftc}), and adaptive control relevant to topology-change resilience (Section~\ref{sec:rw_adaptive}). It synthesizes the literature to identify the gap: no prior work addresses cable-fault tolerance in decentralized cooperative aerial transport or formalizes the mechanism by which N-agnostic control provides implicit fault tolerance under topology changes.
% ============================================================
\subsection{Cooperative Cable-Suspended Aerial Transport}
\label{sec:rw_transport}
% ============================================================

The geometric control foundation begins with Lee, Leok, and McClamroch~\cite{lee2010geometric}, who developed singularity-free tracking on $\SOthree$. Sreenath, Lee, and Kumar~\cite{sreenath2013geometric} extended this to cooperative cable-suspended payload manipulation on $\SEthree \times (\Sph^2)^N$, Lee~\cite{lee2018geometric} proved almost-global exponential stability on $\SEthree \times \Sph^2$ for cable-suspended rigid-body transport, and Dai, Lee, and Bernstein~\cite{wu2014geometric} introduced adaptive load-mass estimation.

Recent advances include agile cooperative aerial manipulation using INDI-based control and centralized NMPC~\cite{sun2025agile}, efficient QP-based cable-force allocation for embedded microcontrollers~\cite{wahba2024efficient}, the validated flexible simulator RotorTM~\cite{li2024rotortm}, and a recent survey of suspended-load transport~\cite{estevez2024review}.

Every published controller assumes intact cables. No result addresses tracking guarantees under cable loss, provides design criteria for sizing against cable-loss disturbances, or characterizes fault tolerance in this system class.

% ============================================================
\subsection{Fault-Tolerant Control and Fault Detection for UAVs}
\label{sec:rw_ftc}
% ============================================================

The UAV fault-tolerant control (FTC) literature is extensive for actuator and sensor faults. Representative results include INDI-based flight after complete loss of two opposing rotors~\cite{sun2020incremental}, fault-tolerant neural CBFs for sensor faults~\cite{zhang2024fault}, and recent surveys of UAV fault detection and diagnosis~\cite{adaika2025fdd}.

\emph{Critical observation:} All studied fault modes are actuator faults (rotor degradation/loss) or sensor faults (IMU bias, GPS dropout). Cable faults, which change the mechanical coupling between agents and payload, are absent.

The only paper addressing cable-related failure in multi-UAV transport is Liang~et~al.~\cite{liang2021fault}, who propose dynamic load reassignment with sliding-mode observers. However, this work assumes known failure identity, requires centralized reconfiguration, provides no detection mechanism, and includes no experimental validation, all of which are incompatible with the N-agnostic setting.

Related cable-failure work outside cooperative multi-UAV transport includes cable-driven parallel robots~\cite{raman2023cable}. None addresses cooperative multi-UAV transport or satisfies the N-agnostic constraint.

% ============================================================
\subsection{Adaptive Control and Disturbance Estimation}
\label{sec:rw_adaptive}
% ============================================================

Two adaptive paradigms are most relevant to topology-change resilience. Concurrent learning (CL)~\cite{chowdhary2010concurrent,chowdhary2015concurrent} replays stored regressor-measurement pairs to achieve parameter convergence without persistent excitation, but after abrupt parameter changes its buffer becomes a liability because pre-change data bias the estimate during turnover. $\mathcal{L}_1$ adaptive control~\cite{cao2010l1book,cao2010stability} provides uniform transient guarantees with memoryless piecewise-constant adaptation, making it structurally suited to topology-change scenarios; Wu~et~al.~\cite{wu2025l1quad} demonstrated $\mathcal{L}_1$ augmentation of geometric quadrotor control at $>$200~Hz. Related decentralized adaptive transport work includes fuzzy-logic adaptation~\cite{barawkar2024decentralized} and adaptation to robot departure~\cite{gao2024decentralized}.

Neither CL nor $\mathcal{L}_1$ has been applied to cooperative cable transport, nor has either been analyzed under abrupt topology changes where the system structure, not just parameter values, changes instantaneously.

Control barrier functions (CBFs)~\cite{ames2017control} provide forward invariance via QP-based control modification; Input-to-State Safety (ISSf)~\cite{kolathaya2019input} extends these guarantees to bounded disturbances, and high-order extensions~\cite{xiao2022control} handle relative-degree constraints. Yang~et~al.~\cite{yang2025robust} combined CBFs with disturbance estimators for cable-suspended payload safety. Complementary tools include contraction-theory tracking-error bounds~\cite{manchester2017control}. None of these methods has been applied to post-fault topology changes in multi-vehicle cable transport.

The theoretical foundations for analyzing topology changes include stability of switched systems under average dwell-time~\cite{hespanha1999stability}, ISS under switching~\cite{vu2007iss}, and the hybrid dynamical systems framework~\cite{goebel2012hybrid} for continuous flow interrupted by discrete jumps---the natural setting for cable-snap events.

The closest empirical observations of emergent fault tolerance are Zeng~et~al.~\cite{zeng2025decentralized} and Gao~et~al.~\cite{gao2024decentralized}. Zeng~et~al.~demonstrated implicit resilience to MAV loss in multi-agent RL for cable-suspended transport with zero-shot sim-to-real transfer, but did not formalize the mechanism. Gao~et~al.~showed adaptation to robot departure, but in fully actuated hexarotors undergoing commanded departure rather than spontaneous cable faults.

No work formalizes the mechanism by which N-agnostic control provides implicit cable-fault tolerance. The switched systems literature provides relevant tools but has not been applied to N-agnostic cable-suspended transport with instantaneous, undetected topology changes.

The contribution lies at the intersection of three well-populated but previously disconnected areas:

\begin{enumerate}[label=(\roman*),nosep]
  \item \emph{Cooperative cable-suspended transport}~\cite{lee2010geometric,sreenath2013geometric,lee2018geometric,sun2025agile,estevez2024review}: mature geometric control with no cable-fault handling.
  \item \emph{UAV fault-tolerant control}~\cite{sun2020incremental,adaika2025fdd,zhang2024fault}: extensive actuator/sensor fault work, but cable faults are absent.
  \item \emph{Adaptive decentralized control}~\cite{chowdhary2010concurrent,cao2010l1book,wu2025l1quad,gao2024decentralized}: rich theory not analyzed for topology-change resilience in cable transport.
\end{enumerate}

No prior work provides tracking bounds or design criteria for cable-fault tolerance in decentralized cooperative aerial transport. The present paper addresses this gap.
